% ================== INTRODUCTION GÉNÉRALE ==================
\chapter*{Introduction Générale}
\addcontentsline{toc}{chapter}{Introduction}
\markboth{Introduction Générale}{Introduction Générale}

La qualité des produits agricoles destinés à la transformation industrielle constitue un enjeu majeur pour la valorisation des chaînes de valeur locales. Au Bénin, le soja occupe une place croissante dans l’économie agricole en raison de son importance alimentaire, industrielle et commerciale. Cependant, les opérations post-récolte réalisées en milieu rural restent souvent limitées par l’insuffisance d’équipements adaptés au nettoyage, au séchage, à la pesée et au suivi de la qualité des lots.

Dans ce contexte, les exigences des unités de transformation imposent une meilleure maîtrise des paramètres de qualité, notamment le taux d’humidité, la pureté du grain et la traçabilité des lots. Les méthodes traditionnelles, fondées principalement sur le tri manuel et le séchage au sol, ne permettent pas toujours de garantir une qualité constante avant l’acheminement vers les centres industriels. Cette situation peut entraîner des pertes post-récolte, des refus de lots et une diminution de la valeur marchande du soja.

Le projet SMART-SOJA s’inscrit dans cette problématique. Il vise à proposer une unité mobile, autonome et connectée capable d’assurer certaines opérations de pré-traitement directement sur les sites de collecte. La solution envisagée combine un sous-système mécanique de nettoyage et de séchage, une chaîne électronique de mesure et de commande, une alimentation photovoltaïque autonome et un dispositif de traçabilité numérique basé sur l’Internet des objets.

Ce projet a été développé au cours de notre stage de fin d'études au sein de l'entreprise \textbf{QOTTO Bénin}, spécialisée dans les systèmes d'énergies renouvelables connectés hors-réseau. Cette immersion nous a offert un cadre favorable à la concrétisation et à la validation expérimentale de ce travail.

La problématique de ce mémoire peut alors être formulée comme suit : comment concevoir une unité mobile et autonome de pré-traitement du soja permettant d’améliorer le nettoyage, le contrôle de l’humidité, la pesée et la traçabilité des lots dans les zones rurales de production au Bénin?

L’objectif général est de \textbf{concevoir} une unité mobile, autonome et connectée de pré-traitement du soja adaptée au nettoyage, au séchage, à la pesée et à la traçabilité numérique des lots en milieu rural béninois.

De manière spécifique, il s’agit de :
\begin{enumerate}
    \item \textbf{Analyser} les besoins de pré-traitement du soja et les exigences de qualité attendues par les unités de transformation;
    \item \textbf{Développer} une solution mécatronique intégrant les sous-systèmes mécanique, électronique, énergétique et logiciel;
    \item \textbf{Évaluer} les performances fonctionnelles du prototype à partir d’indicateurs de nettoyage, d’humidité, de pesée, de consommation énergétique et de traçabilité.
\end{enumerate}

La démarche adoptée repose sur l’analyse du besoin, la définition d’un cahier des charges, la conception fonctionnelle de la solution, le choix des composants, le dimensionnement des principaux sous-systèmes, l’implémentation du contrôle-commande et la réalisation de tests fonctionnels.

Le mémoire est organisé en trois chapitres. Le premier chapitre présente le déroulement du stage, l’organisme d’accueil, les missions confiées et les compétences mobilisées. Le deuxième chapitre est consacré à l’étude de la solution proposée, depuis le contexte de la filière soja jusqu’aux choix mécaniques, électroniques, énergétiques et logiciels. Le troisième chapitre présente les résultats obtenus, leur analyse, leur discussion ainsi que les limites et perspectives du travail.
\newpage
